% 固定假设的集中与数据依赖选择之间的区别。
\begin{tikzpicture}[figure picture,foundation picture,x=1mm,y=1mm,font=\figurefont\footnotesize]
  \draw[foundation flow] (0,0) -- (94,0) node[right] {假设};
  \draw[foundation flow] (0,0) -- (0,48) node[above] {分类错误率};
  \draw[FigureGrid,line width=.5pt] (0,10) -- (92,10);
  \draw[FigureGrid,line width=.5pt] (0,25) -- (92,25);
  \draw[FigureGrid,line width=.5pt] (0,40) -- (92,40);

  \foreach \x/\pop/\emp in {10/31/29,25/22/19,40/27/23,55/18/12,70/24/20,85/20/17}{
    \draw[FigureRed,line width=1.1pt] (\x-3,\pop) -- (\x+3,\pop);
    \fill[FigureBlue] (\x,\emp) circle (1.3);
    \draw[FigureMuted,line width=.5pt] (\x,\emp) -- (\x,\pop);
  }
  \node[font=\figurefont\footnotesize,text=FigureRed,anchor=west] at (5,45) {横线：期望风险$R_{\mathcal P}(h)$};
  \node[font=\figurefont\footnotesize,text=FigureBlue,anchor=west] at (47,45) {圆点：经验风险$\hat R_S(h)$};
  \draw[FigureYellow,line width=1.1pt,rounded corners] (49,7) rectangle (61,34);
  \node[font=\figurefont\footnotesize,text=FigureInk,align=center] at (55,4)
    {ERM输出};
  \draw[foundation flow,FigureYellow] (55,7) -- (55,10.5);
\end{tikzpicture}%
